\section{Discussion and Limitations}
\label{sec:discussion_limitations}

HyDE-style document-like query expansion is associated with improved multi-hop supporting-evidence acquisition under the processed local-corpus protocol studied here. This framing is intentionally empirical: the paper does not claim HyDE itself as a new algorithm, but uses a controlled protocol to test when document-like expansion helps complete supporting-title and paragraph-text retrieval. On both HotpotQA and 2WikiMultihopQA, the observed gain appears first in complete supporting-title retrieval and then in reader EM/F1. Paired retrieval tests, paired answer tests, and supporting-title completeness buckets point in the same direction: examples with all supporting titles retrieved are much easier for the reader than partial-title-support or no-title-support examples. A supporting-paragraph audit over \texttt{is\_supporting=true} corpus records gives identical all-support and recall values for the audited rows listed in Supplement Table~\ref{supp:tab:supporting_paragraph_audit}, reducing the concern that the title metric is inflated by retrieving the right page name but the wrong processed paragraph. Because each annotated supporting title maps to one supporting paragraph record in these released processed splits, the audit is best read as an exact record-recovery check rather than as independent sentence-level evidence. The central empirical message is therefore that supporting-title completeness, aligned with paragraph-text recovery in the processed artifacts, is a major observed bottleneck under the evaluated protocol.

The mechanism evidence is bounded. The matched-call rewrite and equal-budget diagnostics rule out ``one extra LLM call'' and generated length as sufficient explanations, but they do not show that document-like organization universally dominates every high-information query form. Document-like passages improve over direct rewrite and decomposition in the equal-budget diagnostic, while the comparison with keyword/entity expansion is dataset-dependent. The targeted Qwen-Turbo readout gives the same caution downstream: it is positive over direct rewrite but mixed against keyword/entity expansion. We therefore interpret the result as a retrieval-stage advantage of document-like organization over ordinary rewriting and decomposition under this generator, encoder, and local-corpus protocol, not as a reader-independent or universal end-to-end superiority claim.

The guarded canonicalization diagnostic is deliberately secondary: it changes HyDE-style RAG by only +0.0040 EM and +0.0030 F1, is not significant under McNemar's paired EM test ($p=0.5000$), and is reported in the supplement as an answer-format risk analysis rather than as a main contribution.

The main boundary conditions are concentrated in the threats-to-validity analysis below: answer-string overlap in generated query text, the processed local-corpus protocol rather than full-wiki retrieval, and structural rather than latency-measured cost accounting.

\subsection{Threats to Validity}
\label{sec:threats_to_validity}

\textbf{Construct validity.} The evaluation uses supporting-title hit@10, mean supporting-title recall@10, EM, and token F1 as operational measures of evidence acquisition and answer quality. These metrics fit the processed multi-hop QA protocol, but they do not capture all forms of evidence usefulness or answer correctness. The added supporting-paragraph audit tightens the retrieval unit from title to exact normalized title--paragraph text and finds no title-versus-paragraph discrepancy in the audited rows listed in Supplement Table~\ref{supp:tab:supporting_paragraph_audit}; because the released processed splits provide one supporting paragraph record per annotated supporting title, this audit primarily verifies exact processed-record recovery. However, the processed 500-example JSONL artifacts used here do not expose gold supporting sentence text, so we cannot measure supporting-fact sentence coverage or within-paragraph grounding quality. EM and token F1 can also penalize acceptable aliases. The EM-mismatch analysis and qualitative cases make these boundaries visible, but the taxonomy is a single-annotator diagnostic without independent agreement measurement and is not a substitute for systematic human factuality annotation.

\textbf{Internal validity.} The main internal risk is query-side answer-string overlap from the hypothetical passage. This risk is related to recent evidence that LLM-based hypothetical-document expansion can reproduce benchmark knowledge in the generated query text \cite{yoon2025knowledgeleakage}. Many generated hypothetical passages contain the normalized gold answer string: 292/500 on HotpotQA and 347/500 on 2WikiMultihopQA. The reader-only evidence boundary prevents the generated passage from being used as reader evidence, but the passage can still contain answer-like strings and influence which corpus documents are retrieved. Exact answer-absent, gold-content-stratified, and gold-aware scrubbed retrieval diagnostics reduce the explanation that HyDE gains come only from exact normalized answer-string overlap. They also show that answer-like strings and supporting-entity mentions contribute to the full gain. These diagnostics cannot rule out paraphrased answer content, parametric memory, benchmark contamination, or other semantic gold-evidence traces in the query generator.

The token-matched audit truncates existing HyDE passages rather than independently generating equal-length hypothetical passages. It therefore measures sensitivity to information removal rather than providing a clean causal comparison of textual form at equal generation length. The new independent equal-budget diagnostic reduces this limitation for retrieval metrics, but it cannot fully equalize entity count, factual density, or parametric knowledge in generated text, and it does not test reader sensitivity to the retrieved evidence.

A second internal-validity risk is configuration selection and experiment timing. The reported prompt and result artifacts are frozen for audit and rescoring, but not every design choice was selected on a separate development split or before the first HotpotQA result analysis. We therefore distinguish three evidence tiers. The primary frozen study is the original MiniLM-based HotpotQA retrieval-reader comparison with the main lightweight baselines, HyDE query generation, matched-call direct rewriting control, top-$k=10$ evidence budget, and Qwen3.7-Max reader. The BGE-base reproduction, equal-budget 40-word query-form diagnostic, gold-aware scrubbing, corpus-scale and hard-negative audits, top-$k$ sensitivity, supporting-paragraph audit, and Qwen-Turbo fixed-evidence reader check are post-hoc robustness diagnostics added after the primary results to test specific threats to validity. The equal-budget prompts were designed after observing the primary HotpotQA pattern; after the prompt family was fixed, it was applied unchanged to both HotpotQA and 2WikiMultihopQA, making 2Wiki a no-retuning secondary check for that diagnostic rather than a fully independent pre-registered validation. Finally, verifier risk rules and deterministic guards were informed by diagnostics on the reported processed HotpotQA split before final artifact freezing. The verifier-backed canonicalization analysis should therefore be read as an exploratory safety and answer-standardization study, not as evidence of independently tuned verifier generalization.

\textbf{External validity.} The experiments use released 500-example processed HotpotQA and 2WikiMultihopQA splits with local candidate evidence pools. This makes paired comparison and artifact freezing practical, but it is not a full-wiki retrieval setting. The stronger-encoder robustness check, the 50k/100k processed-train distractor stress test, the leakage-filtered and BGE-hard corpus-scale audit, and the top-5/10/20 sensitivity audit reduce the risk that the effect is limited to one weak encoder, the smallest local index, added gold-title leakage, random easy distractors, or a single rank-10 cutoff, but they still do not replace open-domain Wikipedia retrieval. The BGE-hard row is selected from the already audited processed-train 100k pool rather than from all Wikipedia paragraphs, so it tests a harder local negative subset but not the full open-domain retrieval distribution. The implemented comparison rows also have a limited role: BM25, Hybrid, and Dense are lightweight retrieval baselines, while matched-call rewriting, rule-based multi-query retrieval, and rule-based iterative retrieval are mechanism or heuristic controls. They should not be read as faithful public implementations of Query2doc, Rewrite-Retrieve-Read, ReDE-RF, IRCoT, or agentic RAG systems. Results may differ with official or more closely tuned published baselines, other retrievers, different readers, open-weight systems, top-$k$ budgets beyond 20, or datasets with different compositional structure. One targeted Qwen-Turbo readout of fixed equal-budget evidence preserves the document-like answer advantage over direct rewrite on both datasets, while the comparison with keyword/entity expansion is positive on HotpotQA and essentially tied on 2Wiki. This single hosted-reader check is not a broad reader sweep and does not change the paper's retrieval-stage mechanism claim.

\textbf{Reproducibility and cost validity.} Prompts, retrieval files, answer files, analysis scripts, and model invocation ledgers are frozen as project artifacts, but hosted LLM APIs can change over time. The Qwen3.7-Max records document interface, dates, decoding settings, and artifact paths. They do not include provider-side request IDs, exact token accounting for every call, wall-clock latency, or billing exports. The efficiency claim is therefore limited to structural LLM-call counts, retrieval-query executions, and prompt-character counts.
